\setchapterpreamble[u]{\margintoc}
\chapter{\Ac{SIM} Reconstruction}
\labch{sim-reconstruction}

% FIX: Add note about notation... Small letters are spatial domain and capitals are in frequency domain <18-12-23> 

It was laid-out in the theory part, how to achieve introducing higher frequencies in the output image, which would be 
    otherwise lacking in the regular wide-field microscope acquisition outputs.
The process of reconstruction concerns with 
\begin{enumerate}[label=\roman*.]
    \item determining the powers of these aliased frequencies and 
    \item mapping their positions so that they are at the correct corresponding locations in the reconstructed image.
\end{enumerate}
If everything is done correctly, it is possible to recover all the additional information from the newly introduced 
    frequencies in the estimate of the source signal.

There is a multitude of approaches for performing these tasks, each of which has its own benefits and drawbacks.
It is possible to categorise them into three broad categories based on the underlining theory that is used, namely:
\begin{enumerate}[label=(\alph*)]
    \item Wiener filter based,
    \item inversion/optimization based and
    \item machine learning based.
\end{enumerate}

At the input of our reconstruction are 9 \acs{LR}\sidenote{
    In the following text, sensed images (i.e. signal that we measure at the sensor) will be called a \acf{LR} image and
        the reconstructed image will be called a \ac{HR} image. 
    It is important to remember throughout the text that when talking about the sampled (discrete) \ac{LR} image, it has 
        different pixel dimensions (sampling rate) from the \ac{HR} image.
    } images acquired with 3 orientations (phase-shifts) that are spaced by approximately \(2\pi/3\) to maximize 
    information gain in the reconstructed image and every orientation is acquired with 3 phase-offsets 
    (\reffigshort{beads-LR-9-images}).
The acquisition is made in an orientation-major fashion, meaning that all \ac{LR} images of the same orientation are
    sensed before changing to the next orientation.

For most methods from the first two categories it is necessary to estimate some or all of the illumination pattern 
    parameters of the \ac{LR} images prior to the reconstruction.
There again are many possible estimation algorithms relying on various assumptions, which are beneficial for
    acquisitions with specific properties, importantly differing in their robustness to noise contained in the \ac{LR} 
    images.

First, one such estimation sequence for determining the parameters, that has been tested to work well for the data we 
    are working with, will be described (\refsecshort{IP-parameter-estimation}).
The following sections will demonstrate the Wiener filter based approach 
    (\refsecshort{wiener-filter-reconstruction}) and then the inversion/optimization based approach 
    (\refsecshort{inversion-reconstruction}).
Lastly, some of the limitations and benefits of all three approaches will be discussed.

% \begin{marginfigure}
%     \marginpng{otf-model}
%     \caption{The \ac{OTF} model of the optical system.}
%     \labfig{otf-model}
% \end{marginfigure}
\begin{marginfigure}
    \marginpng{otf-model-surf}
    \caption{The \ac{OTF} model of the optical system.}
    \labfig{otf-model}
\end{marginfigure}

In this chapter, we will be using a model \ac{OTF} (\reffigshort{otf-model}) supplied by the assembler of the 
    microscope.
It is derived from the ideal circular aperture \ac{OTF} model with a curvature degradation of \(\alpha = 0.9\).
The formula for the \ac{OTF} is
    \begin{equation*}
        \otf(\vec{k}) = \frac{2}{\pi}\ab(\arccos\ab(\frac{|\vec{k}|}{2k_0}) - \frac{|\vec{k}|}{2k_0}\sqrt{1
        - \frac{|\vec{k}|^2}{4k_0^2}}) \cdot \alpha^{|\vec{k}|/2k_0},
    \end{equation*}
    where \(k_0\) is the cut-off frequency of the \ac{OTF} and for \(\rho\) the pupil radius the emission wavelength 
    \(\lambda_{\text{em}} = \qty{488}{\nano\meter}\), the distance between the entrance pupil and the focal plane \(f_R\) and the 
    numerical aperture of the objective \(\NA\) is
\begin{equation*}
    k_0  = \frac{\rho}{\lambda_{\text{em}} \cdot f_R} = \frac{2\NA}{\lambda_{\text{em}}}.
\end{equation*}

% TODO: Change to margin figure? The differences are not so big anyway <22-11-23> 
\begin{figure}
  \png{beads-LR-imgs}
  \caption{Acquisitions of 9 wide field images with harmonic patterned illuminations. The orientations are labelled
     \(n = 1, 2, 3 \) and the phases \(\phi_\ip = \phi_1,\phi_2,\phi_3\). The order of acquisition is row-major in this
    figure.}
  \labfig{beads-LR-9-images}
\end{figure}%

\section{Parameter Estimation}
\labsec{IP-parameter-estimation}
The general form of the \ac{LR} image illumination patterns in one orientation, denoted \(\ip_{\phi_1}\), 
    \(\ip_{\phi_2}\) and \(\ip_{\phi_3}\), that are at the input of the Wiener filter reconstruction is
    \begin{equation}
        \labeq{ip-formula}
        \ip_{\phi}(\vec{r}) = 1 + \tfrac{m}{2} \cos\ab(2\pi \vec{k}_\ip \cdot \vec{r} + \phi),
    \end{equation}
    where the phases \(\phi_1 \neq \phi_2 \neq \phi_3\) and \(\vec{k}_\ip\) is the wave vector of the illumination pattern.
An assumption that is made in the our acquisition is that the phases differ by \(2\pi/3\), that is 
    \[
        \phi_0 \coloneqq \phi_1 = \phi_2 - 2\pi/3 = \phi_3 - 4\pi/3.
    \]
This assumption is based on the method that is employed for the generation of the illumination patterns, which should
    approximately guarantee this.
Equivalently, due to the \(2\pi\) period of the \(\cos(x)\) function in \refeq{ip-formula}, we adopt the notation
    \begin{gather*}
        \ip_{\phi_1} = \ip_{\phi_0} \eqqcolon \ip_0,\\
        \phi_{-} \coloneqq \phi_0 - 2\pi/3 \quad\&\quad \ip_{\phi_3} = \ip_{\phi_{-}} \eqqcolon \ip_{-},\\
        \phi_{+} \coloneqq \phi_0 + 2\pi/3 \quad\&\quad \ip_{\phi_2} = \ip_{\phi_{+}} \eqqcolon \ip_{+}.
    \end{gather*}
As was demonstrated in the theory part, the acquisition carried out with an illumination pattern \(\ip_\phi\) leads to
    the sensed \ac{LR} image in the Fourier domain \(\Fsensed_\phi \coloneqq \F*{\sensed_\phi}\) being
    \begin{equation}
        \labeq{component-separation}
        \begin{split}
            \Fsensed_\phi &= \F*{\ip_\phi \cdot\source} \cdot \otf
            + \Fnoise  = \ab(\Fip_\phi \conv \Fsource) \cdot \otf + \Fnoise\\[1ex]
                          &= (\underbrace{\ab(\frac{me^{-i \phi}}{4} \dirac_{-\vec{k}_\ip} + \dirac +
                          \frac{me^{i \phi}}{4} \dirac_{+\vec{k}_\ip})}_{\Fip_{\phi}} \conv \Fsource)\cdot\otf + \Fnoise,
        \end{split}
    \end{equation}
    where \(\source\)/\(\Fsource\) is the source signal (i.e. the true structure), \(\otf\) is the \ac{OTF} of the
    imaging system and \(\Fnoise\) is noise in the Fourier domain.
It can be observed from \refeq{component-separation}, that the sensed \ac{LR} image contains 3 bands (components) of 
    aliased signal frequencies with frequency shifts of \(-\vec{k}_\ip\), \(\vec{0}\) and \(+\vec{k}_\ip\).
We can write the three acquisitions in one orientation as
\begin{equation} \labeq{one-orientation-components}
    \begin{pmatrix}[1.2]
        \Fsensed_{\phi_{-}} \\ \Fsensed_{\phi_0} \\ \Fsensed_{\phi_{+}}
    \end{pmatrix} = 
    \begin{pmatrix}[1.2]
        e^{\frac{+2\pi i}{3}}             & 1 & e^{\frac{-2\pi i}{3}} \\
        1 & 1 &  1\\
        e^{\frac{-2\pi i}{3}}             & 1 & e^{\frac{+2\pi i}{3}} \\
    \end{pmatrix}
    \begin{pmatrix}[1.2]
        \frac{m}{4}  e^{-i\phi_0} \cdot \Fcomponent_{-\vec{k}_\ip} \\
        \Fcomponent_{\vec{0}} \\
       \frac{m}{4} e^{+i\phi_0} \cdot\Fcomponent_{+\vec{k}_\ip} \\
   \end{pmatrix}
    + \begin{pmatrix}[1.2]
        N_{\phi_{-}}\\
        N_{\phi_{0}}\\
        N_{\phi_{+}}\\
    \end{pmatrix},
\end{equation}
with the notation
\begin{equation}
    \begin{pmatrix}[1.2]
       \Fcomponent_{-\vec{k}_\ip} \\
       \Fcomponent_{\vec{0}} \\
       \Fcomponent_{+\vec{k}_\ip} \\
   \end{pmatrix}\hspace{-0.3em}(\vec{k}) \coloneqq
   \otf(\vec{k}) \begin{pmatrix}[1.2]
        \dirac_{-\vec{k}_\ip} \conv \Fsource \\
        \Fsource \\
        \dirac_{+\vec{k}_\ip} \conv \Fsource \\
    \end{pmatrix}\hspace{-0.3em}(\vec{k}) =
     \otf(\vec{k}) \begin{pmatrix}
        \Fsource(\vec{k} + \vec{k}_\ip) \\
        \Fsource(\vec{k})  \\
        \Fsource(\vec{k} - \vec{k}_\ip)  \\
    \end{pmatrix}.
\end{equation}

\subsection{Phase-shift Estimate}
\labsubsec{shift-estimation}
First of the parameters that must be estimated is the phase-shift (i.e. wave vector of the illumination pattern) 
    \(\vec{k}_i\).
First a separation of the shifted bands is performed while keeping the multiplicative coefficients\sidenote{The 
    reference phase-offset \(\phi_0\) is yet to be determined.} \(\frac{m}{4} e^{-i\phi_0}\) and 
    \(\frac{m}{4} e^{+i\phi_0}\).
This can be done by solving the linear system in \refeq{one-orientation-components} for 
    \begin{equation}
        \labeq{tilde-components}
        \begin{pmatrix}[1.2]
            \tilde{\Fcomponent}_{-\vec{k}_\ip} \\
            \Fcomponent_{\vec{0}} \\
            \tilde{\Fcomponent}_{+\vec{k}_\ip} \\
        \end{pmatrix} \coloneqq
        \begin{pmatrix}[1.2]
            \frac{m}{4}  e^{-i\phi_0} \cdot \Fcomponent_{-\vec{k}_\ip} \\
            \Fcomponent_{\vec{0}} \\
        \frac{m}{4} e^{+i\phi_0} \cdot\Fcomponent_{+\vec{k}_\ip} \\
    \end{pmatrix},
    \end{equation}
    with disregard for noise.
For the invertible symmetric matrix
    \begin{equation*}
        \tilde{M} \coloneqq \begin{pmatrix}
            e^{\frac{+2\pi i}{3}}             & 1 & e^{\frac{-2\pi i}{3}} \\
            1 & 1 &  1\\
            e^{\frac{-2\pi i}{3}}             & 1 & e^{\frac{+2\pi i}{3}} \\
        \end{pmatrix},
    \end{equation*}
    it is best to find \(\tilde{M}^{-1}\) using Cramer's rule, since many of the minors turn out to be equivalent
    \begin{equation*}
        \begin{split}
            \tilde{M}^{-1} &= 
        \begin{pmatrix}[1.4]
            \phantom{-}\ab|\begin{smallmatrix}
              \tilde{M}_{22} & \tilde{M}_{23} \\ \tilde{M}_{32} & \tilde{M}_{33}
          \end{smallmatrix}| & 
            -\ab|\begin{smallmatrix}
              \tilde{M}_{21} & \tilde{M}_{23} \\ \tilde{M}_{31} & \tilde{M}_{33}
          \end{smallmatrix}| &
          \phantom{-}\ab|\begin{smallmatrix}
              \tilde{M}_{21} & \tilde{M}_{22} \\ \tilde{M}_{31} & \tilde{M}_{32}
            \end{smallmatrix}| \\
            -\ab|\begin{smallmatrix}
              \tilde{M}_{12} & \tilde{M}_{13} \\ \tilde{M}_{32} & \tilde{M}_{33}
          \end{smallmatrix}| & 
          \phantom{-}\ab|\begin{smallmatrix}
              \tilde{M}_{11} & \tilde{M}_{13} \\ \tilde{M}_{31} & \tilde{M}_{33}
          \end{smallmatrix}| &
            -\ab|\begin{smallmatrix}
              \tilde{M}_{11} & \tilde{M}_{12} \\ \tilde{M}_{31} & \tilde{M}_{32}
            \end{smallmatrix}| \\
            \phantom{-}\ab|\begin{smallmatrix}
              \tilde{M}_{12} & \tilde{M}_{13} \\ \tilde{M}_{22} & \tilde{M}_{23}
          \end{smallmatrix}| & 
            -\ab|\begin{smallmatrix}
              \tilde{M}_{11} & \tilde{M}_{13} \\ \tilde{M}_{21} & \tilde{M}_{23}
          \end{smallmatrix}| &
          \phantom{-}\ab|\begin{smallmatrix}
              \tilde{M}_{11} & \tilde{M}_{12} \\ \tilde{M}_{21} & \tilde{M}_{22}
            \end{smallmatrix}|
        \end{pmatrix}^{\transpose} \cdot \frac{1}{|\tilde{M}|}\\[1em]
        &= \begin{pmatrix}
            e^{\frac{+2\pi i}{3}}  & -2\sin\ab(\frac{+2\pi}{3}) & -e^{\frac{-2\pi i}{3}} \\
            e^{\frac{+2\pi i}{3}}  & 2\sin\ab(\frac{+2\pi}{3})   & e^{\frac{+2\pi i}{3}} \\
            -e^{\frac{-2\pi i}{3}} & -2\sin\ab(\frac{+2\pi}{3}) & e^{\frac{+2\pi i}{3}}
        \end{pmatrix}^{\transpose} \cdot \frac{1}{6\sin(\frac{-2\pi i}{3})}  \\[1em]
        &= \begin{pmatrix}
            \frac{-1}{4 \sqrt{3}} + \frac{i}{4} & \frac{-1}{4 \sqrt{3}} + \frac{i}{4} & \frac{-1}{4 \sqrt{3}}
            - \frac{i}{4}\\
            -\frac{1}{2} & \frac{1}{2} & -\frac{1}{2} \\
            \frac{-1}{4 \sqrt{3}} - \frac{i}{4} & \frac{-1}{4 \sqrt{3}} + \frac{i}{4} & \frac{-1}{4 \sqrt{3}}
            + \frac{i}{4}
        \end{pmatrix}.
        \end{split} 
    \end{equation*}

Finally, we have 
    \begin{equation}
        \begin{pmatrix}[1.2]
            \tilde{\Fcomponent}_{-\vec{k}_\ip} \\
            \Fcomponent_{\vec{0}} \\
            \tilde{\Fcomponent}_{+\vec{k}_\ip} \\
        \end{pmatrix}\hspace{-0.3em}(\vec{k}) \approx \tilde{M}^{-1} \cdot 
        \begin{pmatrix}
        \Fsensed_{\phi_{-}} \\ 
        \Fsensed_{\phi_0} \\
        \Fsensed_{\phi_{+}}
        \end{pmatrix}\hspace{-0.3em}(\vec{k}).
    \end{equation}

It is assumed that there is some overlap of the components \(\Fcomponent_{\vec{0}}\) with both
\(\Fcomponent_{-\vec{k}_\ip}\) and \(\Fcomponent_{+\vec{k}_\ip}\) within the support of the \ac{OTF}\sidenote{
    Unfortunately, there is an inverse relationship between how large the overlapping area is (therefore how robust
    the estimation is against noise) and the maximum achievable resolution, because, the larger the norm of
    \(\vec{k}_{\ip}\) is, the higher the maximum aliased frequency of the component, but the lower the overlap with the
    central component \(\Fcomponent_{\vec{0}}\). 
    Ultimately, most acquisitions using a shift estimation based on the cross-correlation make a compromise of these two
    effects and the phase-shift is placed approximately at \(0.95\) of the \ac{OTF} cut-off frequency.
}.
Accordingly, estimate of \(\vec{k}_\ip\) is performed using cross-correlation\sidecite{huang2018}
    \begin{equation*}
        \begin{split}
            \correlation_{+\vec{k}_\ip}(\vec{p}) &= \sum_{\vec{k}} \ab(\otf(\vec{k} + \vec{p})
            \Fcomponent_{\vec{0}}(\vec{k}))\cdot\ab(\otf(\vec{k})\tilde{\Fcomponent}_{+\vec{k}_{\ip}}(\vec{k}
            + \vec{p})) \\[1.2em]
                                 &\hspace{0.3em}\begin{aligned}
                                     =\frac{m}{4}e^{+\phi_0} \sum_{\vec{k}} &\otf(\vec{k}
                                     + \vec{p})\otf(\vec{k})\Fsource(\vec{k}) \\[-0.7em]
                                                                            &  \hspace{1em} \cdot\otf(\vec{k})\otf(\vec{k} + \vec{p})\Fsource(\vec{k} - \vec{k}_\ip + \vec{p})
                                 \end{aligned} \\[1.2em]
                                 &= \frac{m}{4}e^{+\phi_0} \sum_{\vec{k}} \otf^2(\vec{k} + \vec{p}) \otf^2(\vec{k}) \Fsource(\vec{k})
                                 \Fsource(\vec{k} - \vec{k}_\ip + \vec{p}).
        \end{split}
    \end{equation*}
To make the estimate more robust to noise, we can instead use both components in the cross-correlation
    \begin{equation}
        \labeq{cross-correlation-shift}
            \begin{split}
                \correlation(\vec{p}) = &\sum_{\vec{k}} \ab(\otf(\vec{k} + \vec{p})
                \Fcomponent_{\vec{0}}(\vec{k}))\cdot\ab(\otf(\vec{k})\tilde{\Fcomponent}_{+\vec{k}_{\ip}}(\vec{k}
                + \vec{p})) \\
                                                    + &\sum_{\vec{k}} \ab(\otf(\vec{k} - \vec{p})
                \Fcomponent_{\vec{0}}(\vec{k}))\cdot\ab(\otf(\vec{k})\tilde{\Fcomponent}_{-\vec{k}_{\ip}}(\vec{k}
                - \vec{p})).
            \end{split}
    \end{equation}
We see from \refeq{cross-correlation-shift} that \(\mathcal{C}\) is maximized, when the source signal arguments match, 
    \begin{equation*}
        \Fsource(\vec{k}) = \Fsource(\vec{k} - \vec{k}_\ip + \vec{p}) \implies \vec{k}_\ip = \vec{p},
    \end{equation*}
    and consequently the estimate can be made by maximizing the cross-correlation over values of \(\vec{p}\)
    \begin{equation}
        \hat{\vec{k}_{\ip}} = \argmax_{\vec{p}} \mathcal{C}(\vec{p}).
    \end{equation}

In practice, we cannot evaluate the functions \(\Fcomponent_{-\vec{k}_\ip}(\vec{k})\) and 
    \(\Fcomponent_{+\vec{k}_\ip}(\vec{k})\) directly for an arbitrary value of \(\vec{q} = \vec{k} + \vec{p}\)
    / \(\vec{q} = \vec{k} - \vec{p}\), since they are computed from
    the sampled, \ac{LR} images.
The solution to this is to evaluate them for the closest sampled frequency\sidenote{The closest 
    reciprocal of a multiple of the \ac{LR} image sampling step (i.e. pixel size).} \([\vec{q}]\) and performing a shift
    by \(\ab(\vec{q} - [\vec{q}])\) using the Fourier shift theorem.

\subsection{Phase-Offset and Modulation Estimate}
\labsubsec{phase-offset-estimate}

If we expand the components \(\tilde{\Fcomponent}_{-\vec{k}_\ip}\) and \(\Fcomponent_{\vec{0}}\) using 
    \refeq{one-orientation-components} and \refeq{tilde-components} we get
    \begin{align*}
        \tilde{\Fcomponent}_{-\vec{k}_\ip}(\vec{k}) &= \frac{m e^{-i\phi_0}}{4} \Fcomponent_{-\vec{k}_\ip} 
            = \frac{m e^{-i\phi_0}}{4} \otf(\vec{k}) \Fsource(\vec{k} + \vec{k}_\ip), \\
        \Fcomponent_{\vec{0}}(\vec{k}) &= \otf(\vec{k}) \Fsource(\vec{k}).
    \end{align*}
We see that they differ only by the phase-shift \(\vec{k}_\ip\) and the complex coefficient \(me^{-i\phi_0}/4\).
Using the phase-shift estimate \(\hat{\vec{k}}_\ip \eqqcolon \vec{k}_\ip\), estimates of the source signal
    \(\Fsource(\vec{k})\) and can be made from both components
    \begin{align*}
        \tilde{\hat{\Fsource}_{-\vec{k}_\ip}}(\vec{k}) &= \frac{m e^{-i\phi_0}}{4}\hat{\Fsource}_{-\vec{k}_\ip}(\vec{k})
        = \frac{m e^{-i\phi_0}}{4} \ab(\frac{\tilde{\Fcomponent}_{-\vec{k}_\ip}(\vec{k} - \vec{k}_\ip)}{\otf(\vec{k} - \vec{k}_\ip)}), \\
        \hat{\Fsource}_{\vec{0}}(\vec{k}) &= \frac{\Fcomponent_{\vec{0}}(\vec{k})}{\otf(\vec{k})}.
    \end{align*}
Using complex linear regression from values in the intersection of \(\supp\otf(\vec{k} - \vec{k}_\ip)\) and
\(\supp\otf(\vec{k})\), it is then possible to estimate \(\tilde{m}/4 \coloneqq me^{-i\phi_0}/4\) as
    \begin{gather*}
        \Omega = \supp\otf(\vec{k} - \vec{k}_\ip) \cap \supp\otf(\vec{k}) \\
        \hat{\tilde{m}} =  \frac{4 \cdot\sum_{\vec{k} \in \Omega}
            \ab(\tilde{\hat{\Fsource}}_{-\vec{k}_\ip}(\vec{k}))^{*} \cdot \hat{\Fsource}_{\vec{0}}(\vec{k})}{\sum_{\vec{k} \in \Omega}
            \hat{\Fsource}_{\vec{0}}^{2}(\vec{k})}.
    \end{gather*}
The modulation factor and the phase-offset estimates are consequently
\begin{equation*}
    \hat{m} = |\hat{\tilde{m}}| \quad \& \quad \hat{\phi}_0 = \arg(\hat{\tilde{m}}).
\end{equation*}

% \subsection{Alternative Phase-Offset Estimate}%
% \labsubsec{alternative_phase-offset_estimate}
% TODO: Add the other estimation algorithm. This is not the one used for the estimation of modulation and phase-offset <17-12-23> 

\section{Wiener filter Reconstruction}
\labsec{wiener-filter-reconstruction}

The Wiener filter approach consists of:
\begin{enumerate}[label=\roman*.]
    \item \labitm{component-separation} separation of the aliased components  
        \begin{equation*}
            \ab(\Fsensed_{\phi^{n}_{-}}^{n}, \Fsensed_{\phi^{n}_{0}}^{n}, \Fsensed_{\phi^{n}_{+}}^{n})_{n \in \oneto{3}} \mapsto
            \ab(\Fcomponent_{-\vec{k}^{n}_\ip}^n, \Fcomponent_{\vec{0}}^n, \Fcomponent_{+\vec{k}^{n}_\ip}^n)_{n \in \oneto{3}}
        \end{equation*}
    \item \labitm{component-shift} shifting the components to their corresponding locations within the \ac{HR} image\sidenote{
        In practice the components, when shifted, are placed within the \ac{HR} image's grid (in the Fourier domain).
        We are working with a sensor that samples \(512\times 512\) pixels so these are the dimensions of the \ac{LR} 
            sensed images that are due to the imaging equipment. 
        The dimensions of the output \ac{HR} images are not predetermined or restricted by any external constraints and 
            instead are a choice that is necessary to make.
        There is no reason to make the dimensions unnecessarily high since the amount of information in the 
            reconstructed image is constant for any choice that is greater than the Shannon sampling rate of the 
            reconstructed \(\reconstructed(\vec{r})\).
        Assuming that the \ac{LR} images are already sensed with a sampling rate above the Shannon threshold, we can set
            the \ac{HR} image sampling rate to twice that of the \ac{LR} images since the maximum theoretical 
            enhancement in physical resolution (encapsulated information gain) for a reconstruction with linear \ac{SIM}
            is \(2\times\). With this in mind, the target dimensions of the reconstruction that we use is 
            \(1024\times 1024\).
        } 
        \begin{equation*}
            \ab(\Fcomponent_{-\vec{k}^{n}_\ip}^n, \Fcomponent_{\vec{0}}^n, \Fcomponent_{+\vec{k}^{n}_\ip}^n)_{n \in \oneto{3}} \mapsto 
            \ab(\Fcomponent_{-\vec{k}^{n}_\ip}^{n'}, \Fcomponent_{\vec{0}}^{n'}, \Fcomponent_{+\vec{k}^{n}_\ip}^{n'})_{n \in \oneto{3}}
        \end{equation*}
    \item \labitm{otf-shift} determining the \ac{OTF} of the mapped components and the \ac{HR} image 
        \begin{equation*}
            \ab(\Fcomponent_{-\vec{k}^{n}_\ip}^{n'}, \Fcomponent_{\vec{0}}^{n'}, \Fcomponent_{+\vec{k}^{n}_\ip}^{n'})_{n \in \oneto{3}}
            \mapsto \ab(\otf^{n'}_{-\vec{k}^{n}_\ip}, \otf^{n'}_{\vec{0}}, \otf^{n'}_{+\vec{k}^{n}_\ip})_{n \in \oneto{3}}
        \end{equation*}
    \item \labitm{wiener-filter} application of the Wiener filter to the mapped components 
        \begin{equation*}
            \ab(\Fcomponent_{-\vec{k}^{n}_\ip}^{n'}, \Fcomponent_{\vec{0}}^{n'}, \Fcomponent_{+\vec{k}^{n}_\ip}^{n'})_{n \in
           \oneto{3}},
           \ab(\otf^{n'}_{-\vec{k}^{n}_\ip}, \otf^{n'}_{\vec{0}}, \otf^{n'}_{+\vec{k}^{n}_\ip})_{n \in \oneto{3}}
            \mapsto \Freconstructed.
        \end{equation*}
\end{enumerate}

\subsection{Separation of Components (\refitmshort{component-separation})}
\labsubsec{component-separation}

The procedure for \refitm{component-separation} was already described in \refsubsec{shift-estimation}.
\marginnote{The necessity of the separation of components is one of the major drawbacks for the Wiener filter approach,
    since it requires the acquisition of 9 \ac{LR} images. We used all of the 9 \ac{LR} images for the parameter
    estimation (\refsecshort{IP-parameter-estimation}), but if the parameters where known as a consequence to the method
    of acquisition, or they where estimated using a lower number of \ac{LR} images, it wouldn't be possible to perform
the reconstruction using the Wiener filter approach. It would however be possible to use the other two methods.}
At this moment, however, we have the estimates of the phase-offsets \(\phi_{-}, \phi_{0}, \phi_{+}\) and the 
    modulations \(m\),  therefore, we use them to separate the weighted components \(\Fcomponent_{-\vec{k}_\ip}\), 
    \(\Fcomponent_{\vec{0}}\) and \(\Fcomponent_{\vec{k}_\ip}\).
The linear combination matrix is now
    \begin{equation*}
        \begin{split}
            M(m, \phi_0) &= \begin{pmatrix}[1.2]
            \frac{m}{4} e^{\frac{+2\pi i}{3} - i\phi_0}             & 1 & \frac{m}{4} e^{\frac{-2\pi i}{3}+i\phi_0} \\
            \frac{m}{4}  e^{-i\phi_0} & 1 &  \frac{m}{4} e^{+i\phi_0}\\
             \frac{m}{4} e^{\frac{-2\pi i}{3} - i\phi_0}            & 1 &  \frac{m}{4} e^{\frac{+2\pi i}{3}+i\phi_0}\\
        \end{pmatrix} \\
        &= \tilde{M}\cdot \pdiagmat[ empty = {\cdot} ]{\frac{m e^{ -i\phi_0}}{4}, 1,  \frac{me^{+i\phi_0}}{4} },
        \end{split}
    \end{equation*}
    and \(M^{-1}\) is
    \begin{equation}
        \begin{split}
            M^{-1}(m, \phi_0) &= \pdiagmat[ empty = {\cdot} ]{\frac{4 e^{ +i\phi_0}}{m}, 1,  \frac{4e^{-i\phi_0}}{m} } \cdot
            \tilde{M}^{-1} \\
                &= \begin{pmatrix}[1.2]
                (\frac{-1}{\sqrt{3}} + i) \cdot \frac{e^{+i\phi_0}}{m} & \frac{-1}{4 \sqrt{3}} + \frac{i}{4} & (\frac{-1}{\sqrt{3}}
                - i) \cdot \frac{e^{-i\phi_0}}{m}\\
                -\frac{2e^{+i\phi_0}}{m} & \frac{1}{2} & -\frac{2e^{-i\phi_0}}{m} \\
                (\frac{-1}{\sqrt{3}} - i) \cdot \frac{e^{+i\phi_0}}{m} & \frac{-1}{4 \sqrt{3}} + \frac{i}{4} & (\frac{-1}{\sqrt{3}}
                + i) \cdot \frac{e^{-i\phi_0}}{m}
            \end{pmatrix}.
        \end{split}
    \end{equation}
For every orientation \(n \in \oneto{3}\) the separated components are
    \begin{equation*}
        \begin{pmatrix}
        \Fcomponent^{n}_{-\vec{k}^{n}_\ip} \\
        \Fcomponent^{n}_{\vec{0}} \\
        \Fcomponent^{n}_{+\vec{k}^{n}_\ip} \\
    \end{pmatrix}\hspace{-0.3em}(\vec{k}) 
        \approx M^{-1}(m^{n}, \phi^{n}_0) \cdot \begin{pmatrix}
            \Fsensed^{n}_{\phi^{n}_{-}} \\ \Fsensed^{n}_{\phi^{n}_0} \\ \Fsensed^{n}_{\phi^{n}_{+}}
        \end{pmatrix}\hspace{-0.3em}(\vec{k}).
    \end{equation*}

\begin{figure}
    \png{separated-components}
    \caption{Separated components for orientation \(n=1\).}
    \labfig{separated-components}
\end{figure}

\subsection{Shifting Components (\refitmshort{component-shift} and \refitmshort{otf-shift})}
\labsubsec{shifting-component}

It is necessary need to shift the components to their locations corresponding to their location within the
    source signal.
For a component \(\Fcomponent_{\vec{p}}\) which corresponds to the element of the \(\Fip_{\vec{p}}\) with 
    \(\dirac_{\vec{p}}\) we want to get \(\Fcomponent_{\vec{p}}'(\vec{k}) = \Fcomponent_{\vec{p}}(\vec{k} + \vec{p})\).
This is done using the Fourier shift theorem as
    \begin{equation*}
        \Fcomponent'_{\vec{p}}(\vec{k}) = \F{\Finverse{\Fcomponent_{\vec{p}}} \cdot e^{-2\pi i  \vec{p} \cdot \vec{r}}}.
    \end{equation*}
In this way, the shifting in \refitm{component-shift} is performed for every component \(\Fcomponent_{-\vec{k}^{n}_\ip}\),
    \(\Fcomponent_{\vec{0}}\) and \(\Fcomponent_{+\vec{k}^{n}_\ip}\) and for every orientation \(n \in \oneto{3}\)
    (\reffigshort{shifted-components}).


\floatsetup[figure]{ % Captions for figueres
    capbesideposition={top,outside},%
}%
\begin{figure}
    \png{shifted-components}
    \caption{The components shifted to their corresponding locations in the \ac{HR} image.}
    \labfig{shifted-components}
\end{figure}
\floatsetup[figure]{ % Captions for figueres
    capbesideposition={bottom,outside},%
}%

It is necessary to realize, how the frequencies in component \(\Fcomponent'_{\vec{p}}\) manifest themselves in the
    sensed \ac{LR} images \(\Fsensed\).
The whole point of illuminating the image in a way that shifts the components towards lower frequencies is, that it
    moves them within the support of the \ac{OTF}.
Accordingly to \refeq{component-separation}, the frequencies of the component \(\Fcomponent'_{\vec{p}}\) are the
    projections of the source signal \(\Fsource(\vec{k} - \vec{p})\) frequencies, as if \(\vec{p}\) was the origin of
    the Fourier transformed sensed signal \(\Fsensed\).
That is, an estimate of the source signal from the information contained in the component \(\Fcomponent'_{\vec{p}}\) can
    be illustratively\sidenote{The estimate is not obtained in this way precisely as will be demonstrated in
    \refsubsec{wiener-filter}.} written as
    \begin{equation}
        \labeq{component-with-OTF-mapping}
        \hat{\Fsource}_{\Fcomponent'_{\vec{p}}}(\vec{k}) \cdot \otf(\vec{k} - \vec{p})
        = \Fcomponent'_{\vec{p}}(\vec{k}).
    \end{equation}
As indicated in \refitm{otf-shift} an \ac{OTF} is assigned\sidenote{
    When using a model \ac{OTF}, this is trivially done by evaluation with translated arguments. If a \ac{LR} measured 
    transfer function is to be used it must be upsampled due to the mismatch of sampling rates of the \ac{HR} and 
    \ac{LR} images, and the translation must be realized using the Fourier shift theorem. Alternatively, the evaluation
    may be done of an interpolation of the measured data.
    } to every component \(\Fcomponent'_{\vec{p}}\) of every orientation \(n \in \oneto{3}\)
    \begin{equation*}
        \otf'_{\vec{p}}(\vec{k}) = \otf(\vec{k} - \vec{p}).
    \end{equation*}

% \begin{figure}
%   \png{shifted-otfs}
%   \caption{\ac{OTF}s \(\otf_{\vec{p}}'\) of the \(\vec{p}\)-shifted components.}
%   \labfig{shifted-otfs}
% \end{figure}

\begin{marginfigure}[*-16]
  \marginpng{shifted-otfs}
  \caption{\ac{OTF}s \(\otf_{\vec{p}}'\) of the \(\vec{p}\)-shifted components.}
  \labfig{shifted-otfs}
\end{marginfigure}

\subsection{Wiener Filter (\refitmshort{wiener-filter})}
\labsubsec{wiener-filter}

The state of the reconstruction is that we have frequency components, positioned at locations within the Fourier space 
    that are in compliance with the physical structure that they represent in the source signal and have appropriate 
    \acs{OTF}s for them.
Each \ac{OTF} corresponding to a component acts as degradation function for the source signal \(\Fsource\) in the
    components location (\refeqshort{component-with-OTF-mapping}).
The degradation is embodied by a convolution of the source signal with a \ac{PSF} or in the frequency space
    by multiplication with the \ac{OTF}.
% TODO: Figure of degraded shifted components without deconvolution <18-12-23> 
Inversion of this degradation process in image processing is called deconvolution (also restoration).
There are many approaches with various purposes that attempt to achieve deconvolution each targeted for a different
    purpose.

One such approach, that is widely used is the \emph{Weiner} deconvolution (filter). 
It is derived in as a least square error estimate of the original signal, given a noisy sensed signal and a known
    transfer (degradation) function.% (\refch{wiener-filter}).
A simplified version of the Weiner filter estimate \(\Freconstructed\) of a noisy signal \(\Fsensed\) and an \ac{OTF} 
    \(\otf\) is given by
    \begin{equation*}
        \Freconstructed = \frac{\otf^*\Fsensed}{\otf^*\otf + \omega^2},
    \end{equation*}
    where \(\omega\) is an empirically chosen real constant, which represents inverse \ac{SNR} in the noisy signal.
For the purpose of restoring the relocated components the estimate is realized using frequency information from all of 
    the components in the numerator and the full \ac{HR} image \ac{OTF} in the denominator
    \begin{equation*}
        \Freconstructed_{\text{\ac{SIM}}}(\vec{k}) = \frac{\sum\limits_{n = 1}^{3}\sum_{\vec{p} \in \{-\vec{k}^{n}_\ip,
        \vec{0}, +\vec{k}^{n}_\ip\}} \ab(\otf^{n'}_{\vec{p}}(\vec{k}))^*
        \Fcomponent^{n'}_{\vec{p}}(\vec{k})}{\sum\limits_{n = 1}^{3}\sum_{\vec{p} \in \{-\vec{k}^{n}_\ip,
        \vec{0}, +\vec{k}^{n}_\ip\}} \ab(\otf^{n'}_{\vec{p}}(\vec{k}))^*
\otf^{n'}_{\vec{p}}(\vec{k}) + \omega^2}.
    \end{equation*}

% TODO: Should be an option of the image <19-12-23> 
% NOTE: This setup is to avoid overlapping of captions of this image with the next image  
\floatsetup[figure]{ % Captions for figueres
    capbesideposition={top,outside},%
}%
\begin{figure}
    \png{weiner-reconstruction}
    \caption{Reconstructed \ac{HR} image using \emph{Weiner} deconvolution.}
    \labfig{weiner-reconstruction}
\end{figure}
\floatsetup[figure]{ % Captions for figueres
    capbesideposition={bottom,outside},%
}%

% TODO: Side note about how an optimal Fourier domain image should look like <19-12-23> 
It is visible (\reffigshort{weiner-reconstruction}), in particular in the Fourier domain, that the
reconstruction is not optimal (\reffigshort{weiner-reconstruction-artifacts}). 
The first aspect that is clearly not expected in an optimal reconstruction are noticeable peaks at the locations of the
    shift frequencies \(\vec{k}_\ip^n\) for all orientations \(n \in \oneto{3}\).
\begin{marginfigure}[*-24]
    \marginpng{weiner-reconstruction-artifacts}
    \caption{
    At the top the full reconstructed image without any modifications and filtering.
    There are visible periodic artefacts, in the \ac{SIM} terminology aptly named \emph{honey comb} artefacts, that are 
        most prominent in the background area of the image. 
    In the middle, the same reconstruction but after performing strong apodization filtering. 
    It is apparent that the apodized image has less artefacts in the background but is also noteworthy, that the image 
        also has a lower resolution that can be observed by comparing the sizes of the beads in the filtered and 
        non-filtered reconstructions. 
    At the bottom, difference between the two reconstructions that accents the form of \emph{honey comb} artefacts.
}
    \labfig{weiner-reconstruction-artifacts}
\end{marginfigure}

There is no physical reason for the presence of these peaks.
One explanation is incorrect parameter estimation that leads to an incorrect component separation.
This is partly the case as we can see that the left-top and right-bottom components have slightly more prominent peaks
    than the other two orientations, which suggests that there are "leftover" parts of the central component, which
    should physically have the strongest intensity within the Fourier domain.


However, this reasoning fails to fully explain the deficiency of the reconstruction which reveals itself by the pronounced
    circumference of the Fourier spectrum.
This cannot be explained by inferior component separation, but instead points toward a suboptimal transfer function
    used for the reconstruction which admits frequencies that should be dismissed as being beyond the \ac{OTF} support 
    of the imaging system.
We will tackle this issue in \refch{tf-measurement}.

\subsection{Apodization}
\labsubsec{apodization}
One more tool which commonly gets employed to counter these high-frequency artefacts is apodization of the output 
    Fourier domain \ac{HR} image.
% TODO: Figure apodization relief <kunzaatko> 
Apodization is in effect a low-pass filter which is constant \(1\) near the origin and beyond some predefined radius, 
    decreases to \(0\) according to some smooth function that is termed the \emph{apodization function}. 
The smoothness of the decrease is important for suppressing the formation of artefacts that occur from using the 
    \ac{DFT} on a signal with abrupt edges.
There are many possibilities to use as apodization functions, but in this particular use-case, there is no significant 
    difference between using any of them.
One that is used commonly is the \emph{cosine apodization function}
    \begin{equation*}
        A_{a}(r) = \cos\ab(\frac{\pi r}{2a}),
    \end{equation*}
    where \(a\) is a chosen constant parameter of the apodization.
\begin{marginfigure}
    \marginpng{weiner-reconstruction-surface-apo}
    \caption{The apodization filter used in \reffig{apo-weiner}.}
    \labfig{apo-surface}
\end{marginfigure}

A disadvantage of using apodization for countering these artefacts is that it dismisses the information in the image
    from the high frequencies and essentially acts as a denoising procedure during which inherently some features of the
    image are lost.
Due to this, correctly determining and eliminating the attributes that cause the formation of artefact in the
    reconstruction in the first place is regarded as a better, yet more demanding and non-universal, approach.
The artefacts that are filtered from the original non-apodized image are shown in 
    \reffig{weiner-reconstruction-artifacts} on a cut-out from the spatial reconstructed image, and the full apodized
    reconstruction is in \reffig{apo-weiner}.

\begin{figure}[h]
    \png{weiner-reconstruction-apodization}
    \caption{Reconstructed image after using apodization to suppress artefacts of incorrect reconstruction parameters
    and \ac{OTF}.}
    \labfig{apo-weiner}
\end{figure}


\section{Inversion Reconstruction}
\labsec{inversion-reconstruction}

The principle of the inversion reconstruction is in many ways easier compared to the Weiner filter reconstruction
    (\refsecshort{wiener-filter-reconstruction}).
It relies on formulating a forward model \(\fwd\) that, when applied to the true source signal \(\source\) yields the
    expectation of the sensed \ac{LR} images \(\hat{\sensed}\).
Unlike in the Wiener filter reconstruction, we do not attempt to develop an estimate for \(\fwd^{-1}\), but instead, it 
    is sufficient to formulate the forward model and rely on mathematical optimization for estimating the correct 
    desired input \(\source\), minimizing some defined loss \(\loss\) between \(\fwd[\source]\) and the sensed 
    \(\sensed\).
A simplified outline of the solution that inversion reconstruction gives is
    \begin{equation*}
        \reconstructed = \argmin_{\source} \loss[\sensed,  \fwd[\source] ].
    \end{equation*}


\subsection{Forward Model}
\labsubsec{forward-model}

The forward model for one \ac{LR} acquisition \(\sensed\) with a specified illumination \(\ip_\sensed\) is apparent from
    \refeq{component-separation}
    \begin{equation}
        \labeq{fwd-one-lr}
        \hat{\sensed} = \ab(\ip_\sensed \cdot\source) \conv \psf,
    \end{equation}
    where \(\psf\) is the \ac{PSF} of the optical system.
One difference that is necessary to make from \refeq{fwd-one-lr}, is, that the optimized \(\reconstructed\) must 
    minimize concurrently the loss \(\loss\) for all of the forward model expectations of the \ac{LR} acquired images at
    the input, with a single estimate of \(\reconstructed\).
An obvious and easiest way to achieve this is to minimize the sum of loss \(\ell\) evaluated for every input image
    \begin{equation*}
        \loss_\Sigma[\{\sensed_n\},  \{\fwd_n[\source]\} ] = \sum_{n} \loss[\sensed_n,  \fwd_n[\source] ].
    \end{equation*}

\subsection{Loss Function}
\labsubsec{loss-function}

A choice that is necessary to make is, the form of the loss \(\loss\).
The role of the loss function is mapping the parameters of the optimization problem (in our case the values of the
    pixels of the forward model carried out on the source signal \(\source\)) onto a real number.
The important aspect of loss that is necessary to redeem is, that minimization of the loss leads to a better estimate of
    the parameters that we want to determine.
One very common choice for the loss in deep learning applications is the \(\ell_2\) norm between the measurement and the
    prediction.
It is possible however, to choose a better loss, that is directly linked to the nature of our optimization problem.
An ideal choice is dependent on the nature\sidenote{Similarly to mathematical statistics, there is no optimal choice 
    for an estimate. Instead we can choose to value some of its properties such
    as it being non-biased, its speed of convergence or its robustness to spurious input data (such as outlying
    measurements). From a standpoint of mathematical optimization, also the numerical stability of the optimization must
    be taken into account.} of the estimate \(\reconstructed\) that we want to achieve.
In microscopy (especially in low \ac{SNR} acquisitions), the dominant source of noise is due to photon (Poisson) shot noise
    (\refsubsecshort{noise}).
Assuming a correct forward model \(\fwd\), we can specify the probability of a sensed signal in relation to the true 
    source signal \(\source\) as
    \begin{equation*}
        \labeq{poisson-probablility-sensed}
        \distribution{P}[\sensed \mid \fwd[\source]] = \prod_{\vec{r}}
        \frac{\ab(\fwd[\source](\vec{r}))^{\sensed(\vec{r})}}{\Gamma(\sensed(\vec{r}) + 1)} \cdot e^{-\fwd[\source](\vec{r})}.
    \end{equation*}
Based on this, assumption, we want to get a reconstruction \(\reconstructed\) such that the probability
    \(\distribution{P}[\sensed(\vec{r}) \mid \fwd[\reconstructed](\vec{r})]\) is as large as possible, since this would
    be the reconstruction that is most probable to describe the sensed \ac{LR} images \(\sensed\).
To reach this goal the loss is set to be the negative log-likelihood\sidenote{
        Using this loss during the optimization, if we disregarded the regularization term (\refsubsecshort{regularizer}), is 
            equivalent to performing the \ac{MLE} of the \(\source\) in mathematical statistics.
    } of the sensed \ac{LR} images (taking \(e\) for the 
    base of the logarithm)
    \begin{equation}
      \labeq{log-likelihood-poisson-loss-long}
      \begin{split}
          \tilde{\ell}_{\distribution{P}}[\sensed, \fwd[\source]] &= -\ln(\distribution{P}[\sensed(\vec{r}) \mid
          \fwd[\source](\vec{r})]) \\[0.7em]
                                                                            &= \sum_{\vec{r}}\fwd[\source](\vec{r}) + \ln(\Gamma(\sensed(\vec{r}) + 1)) \\[-0.7em]
                                                                            &\hspace{3em}- \sensed(\vec{r})\cdot\ln(\fwd[\source](\vec{r})).
      \end{split}
    \end{equation}
The term \(\ln(\Gamma(\sensed(\vec{r}) + 1))\) in the loss of \refeq{log-likelihood-poisson-loss-long} is constant for
    any source \(\source\) so we can ignore it.
The final loss for the optimization problem then becomes
\begin{equation}
      \labeq{log-likelihood-poisson-loss}
          \ell_{\distribution{P}}[\sensed, \fwd[\source]] = \sum_{\vec{r}}\fwd[\source](\vec{r})- \sensed(\vec{r})\cdot\ln(\fwd[\source](\vec{r})).
\end{equation}

\subsection{Regularization Term}
\labsubsec{regularizer}

To prevent over-fitting of the reconstructed \(\reconstructed\) during the optimization, it is necessary to add
    regularization to the optimization problem.
Regularization leads to a simpler solution, that may be more robust to the unwanted influence of noise that would 
    otherwise be inflicted by the optimization.
In general, it is necessary to add some form of regularization into the optimization problem, if the ratio  of the 
    number of parameters in the output to number of input data points is large.
The principle of regularization is to impose some additional constraint to the output parameters that forces against
    outputs that would be in conflict with what is physically expected of the output.


Here the regularization is added explicitly in the form of a regularization term.
The form of the regularization term can highly influence the output and so it is necessary to choose it carefully.
Most regularization terms used in the context of images are defined in favour of the principle, that the output should 
    be reasonably smooth.
In other words, it attempts to penalize variation in the output.

One of the most commonly used regularization terms is \emph{\ac{TV}} in the form
\begin{equation*}
    \reg_{\text{TV}}[\source] = \sum_{\vec{r}}| \nabla \source(\vec{r})|,
\end{equation*}
where \(\nabla\source\) is approximated for the discretely sampled \(\source\) and the resulting formula, for \(u\) and
    \(v\) being the pixel indices, is given by
\begin{equation*}
    \reg_{\text{TV}}[\source] = \sum_{u,v} \sqrt{| \source(u + 1, v) - \source(u,v)|^2 + | \source(u, v + 1)
    - \source(u,v)|^2}.
\end{equation*}

Another option that is frequently used for images is \emph{\ac{GR}} for which the continuous formula is
\begin{equation*}
    \reg_{\text{GR}}[\source] = \sum_{\vec{r}} \sqrt{\source(\vec{r})}(\Delta \sqrt{\source})(\vec{r}),
\end{equation*}
where \(\Delta \sqrt{\source}\) is the Laplace operator and it is approximated in the discrete domain as
\begin{multline*}
        \reg_{\text{GR}}[\source] = \sum_{u,v} \sqrt{\source(u,v)} \cdot \left(\frac{\sqrt{\source(u + u_{\Delta}, v)}
        + \sqrt{\source(u - u_{\Delta}, v)}}{u_{\Delta} v_{\Delta}} \right. \\
        + \left.\frac{\sqrt{\source(u, v + v_{\Delta})} + \sqrt{\source(u, v - v_{\Delta})}
        - 4\cdot \sqrt{\source(u,v)}}{u_{\Delta} v_{\Delta}}\right).
\end{multline*}

% TODO: Add note about what is u_Δ  and v_Δ <19-12-23> 

\subsection{Optimization}
\labsubsec{optimization}

The final reconstruction optimization problem is posed as
\begin{equation*}
    \reconstructed =  \argmin_{\source} \ab(\loss_{\Sigma}\ab[\{\sensed_n\},\{\fwd_n[\source]\}]
    + \lambda\reg[\source]),
\end{equation*}
    where \(\lambda\) is the regularization parameter, that is set empirically and the loss is the sum of Poisson
    losses \(\loss_{\distribution{P}}\).

The numerical optimization itself is performed using the \ac{LBFGS} algorithm\sidecite{kmogensen2018} similar to 
    gradient descent with the use of source-to-source backward propagated automatic differentiation \sidecite{innes2018} 
    for computing the gradient and Hessian.

\floatsetup[figure]{ % Captions for figueres
    capbesideposition={top,outside},%
}%
\begin{figure}
    \png{optimization-reconstruction-both}
    \caption{Reconstructed \(\reconstructed\) using optimization reconstruction with \ac{TV} (top row) and \ac{GR}
        (bottom row) for regularization and \(\lambda = 0.05\).
        A big difference in the two reconstructions is not visible in this case.}  % TODO: Add more context to why there
        % is not much of a difference <26-12-23> 
    \labfig{reconstruction-optimization-both}
\end{figure}
\floatsetup[figure]{ % Captions for figueres
    capbesideposition={bottom,outside},%
}%

% TODO: Either the figures must be wide over the full page along with the weiner reconstruction or there must be some
% text in-between the figures or the image must be in one image <26-12-23> 

% \begin{figure}
%     \png{optimization-reconstruction-GR_0_05}
%     \caption{Reconstructed \(\reconstructed\) using optimization reconstruction with \ac{GR} for regularization
%         and \(\lambda = 0.05\).} 
%     \labfig{reconstruction-GR_0_05}
% \end{figure}

\begin{marginfigure}[*-17]
    \marginpng{optimization-reconstruction-compare-regularization}
    \caption{Comparison of the two reconstructions using \ac{GR} and \ac{TV} for regularization (\(\lambda = 0.05\)) on
    a cut-out of the image. It is important to note that \(\lambda\) does not have the same effect for one value on any
    regularization.}
    \labfig{comparison-GR-TV}
\end{marginfigure}

\section{Discussion}%
\labsec{reconstruction_methods_discussion}

There has been a lot of research on the use of different reconstruction methods and their benefits have been
    demonstrated on particular datasets.
% TODO: Add reference to comparison paper for reconstruction methods <26-12-23> 
To this date, no reconstruction method has been distinguished as the "holy grail" that is best for any dataset
    universally, as is often the case for any set of algorithms developed for a particular purpose.
Benefits can be attributed to each method.
More limiting are the drawbacks that may deem an approach utterly inappropriate for a particular methodology of 
    observation.

One such example is the use of machine learning based methods.
The major drawback that they inherently posses is the lack of fidelity of the output.
This is due to the way machine learning methods work in most implementations.
Usually the training dataset is generated by modifying the source similarly to how it was expressed in 
    \refsubsec{forward-model} and the user formulating a procedure to generate model sources \(\source\) that are
    applicable for the given observation.
Generating the training dataset in this manner ensures that the learning can be supplied with a ground truth of the
    source signal.
This means that the training dataset is not necessarily natural in the sense that the data used for training may not be
    variable enough to encompass the structure of the real world data input for reconstruction and the resulting
    model can be biased in the direction of data that it was supplied with during training.
The quality of the model that is problematic in this regard is termed in literature as the \emph{ability of the model to
    generalize} and although it is possible to minimize the issue to a certain extent, it is always present at some 
    degree.
Due to this, machine learning methods are hard to justify in the reconstruction of biological structures which are not 
    yet fully understood to the level of being able to generate models of their structure that are absolutely 
    dependable.
On the other hand, the fact that the reconstruction is "aware" of the source structure that it should output can lead to
    machine learning reconstructions, that use appropriate source models, to be vastly less noisy and visually and 
    structurally more appealing than reconstruction based on other methods.

Another benefit is that the reconstruction using machine learning can be supplied with data without any preparation
    steps.
Denoising and other procedures that ensure the visual quality of the reconstruction can be incorporated into the
    reconstruction step itself which means that there is no information lost before the actual reconstruction is 
    performed.
In general, there is a trend for new algorithms to consolidate all the steps into one single model precisely due to this
    reasoning.

An aspect that must be considered when choosing a reconstruction method, if the object of the observation is not static 
    (e.g. \emph{in vivo} acquisitions), is the value of temporal resolution and the photon efficiency of the 
    observation.
There is an intrinsic deficiency of the Wiener filter reconstruction in the fact that it is necessary to sense the image
    \(3\) times in a single orientation to be able to linearly separate the components.
Machine learning and optimization based reconstruction procedures do not posses this limitation, as long as we are able 
    to estimate the parameters of the illumination patterns.
This can be recognized from the \refeq{ip-formula}, which suggests that we gain no additional information from the two
    additional acquisitions having different phase-offsets.
An equivalent optimization reconstruction is therefore in theory possible to perform using only \(3\) \ac{LR} images, 
    one for each orientation (\reffigshort{optimization-reconstruction-3LR-vs-9LR}).
The difference in these two reconstructions should be explainable only through the additional optimization parameters
    that should benefit the robustness towards noise.

\begin{marginfigure}[*-18]
    \marginpng{optimization-reconstruction-3LR-vs-9LR}
    \caption{Comparison of optimization reconstruction performed from 3 and 9 \ac{LR} images. }
    \labfig{optimization-reconstruction-3LR-vs-9LR}
\end{marginfigure}

A further, very interesting, possible enhancement to the machine learning and optimization approaches is to apprehend
    the estimated parameters of the illumination pattern as random variables with a certain distribution instead of 
    constant estimates or to perform the estimates as part of the reconstruction.
That can be leveraged in the reconstruction to further fine tune the properties we are expecting from the output
    without overly relying on the correctness of the estimates.
Within the machine learning approach this can be achieved by training the reconstruction model without supplying the 
    correct parameters as part of the data.
For the optimization reconstruction, this can be done by Bayesian inversion\sidecite{orieux2012} using Bayes rule in the
    form
    \begin{equation*}
        \distribution{P}[\source, \lambda \mid \sensed] = \frac{\distribution{P}[\sensed \mid \source,
        \lambda] \cdot \distribution{P}[\source, \lambda]}{\distribution{P}[\sensed]},
    \end{equation*}
    where \(\lambda\) are the parameters of the illumination patterns during acquisition of \(\sensed\)
    \(\distribution{P}[\source , \lambda] = \distribution{P}[\source] \cdot \distribution{P}[\lambda]\) due to
    independence, \(\distribution{P}[\source]\) can be set either as the uninformative prior or as a conjugate prior
    to the noise distribution \(\distribution{P}[\sensed \mid \source, \lambda]\) and the prior 
    \(\distribution{P}[\lambda]\) can be formed by randomizing the parameter estimates.

One hybrid approach to integration of parameter estimation into the reconstruction is tiled \ac{SIM}
    reconstruction.
The idea is that the \ac{LR} images are subdivided into tiles and the reconstruction is performed for each tile
    individually.
The parameters are estimated independently for the tiles.
This has been shown to lead to better results in some cases\sidecite{hoffman2020}.

% TODO: Note about tiled reconstruction <26-12-23> 
    
% TODO: The reconstruction should be shown of the CCPs instead of the beads <26-12-23> 
